%%%%%%%%%%%%%%%%%%
\section{Anosov representations in $SL_3\C$}

Consider the collection of Anosov representations 
\begin{align}
    \mathcal{A}_\R(\Gamma) := \{\rho:\Gamma \rightarrow SL_3\R \text{ Anosov}\}\text{, and }
    \mathcal{A}(\Gamma) := \{\rho:\Gamma \rightarrow SL_3\C \text{ Anosov}\}.
\end{align}
Let $\rho_0 : \Gamma \rightarrow PSL_2\R$ Fuchsian and 
$\rho_i : \Gamma \rightarrow PSL_2\C$ quasi-Fuchsian. 
To create Anosov representations in $SL_3\C$, 
we can compose with an irreducible representation $\iota : PSL_2\R \rightarrow SL_3\R$:
\begin{align}
    \rho_{F} = \iota \circ \rho_0 : \Gamma \rightarrow SL_3\R\text{, and }
    \rho_{QF} = \iota \circ \rho_i : \Gamma \rightarrow SL_3\C.
\end{align}
Similarly, (assuming that $\rho_0,\rho_i$ lift to $SL_2\R$) 
we can compose with a reducible representation 
$\Delta : SL_2\C \rightarrow SL_3\C$:
\begin{align}
    \rho_{R} = \Delta \circ \rho_0 : \Gamma \rightarrow SL_3\R\text{, and }
    \rho_{QR} = \Delta \circ \rho_i : \Gamma \rightarrow SL_3\C.
\end{align}
The representations $\rho_{Fuchs}$ and $\rho_{Red}$ are Anosov. 
We can then define the following subsets of $\mathcal{A}$:
\begin{itemize}
    \item $SL_3\C$ quasi-Fuchsian representations
    $\mathcal{H}(\Gamma) := \{\rho \in \mathcal{A}(\Gamma) : \rho = \iota \circ  \text{in the same Anosov component as }\rho_0\}$
    \item Hitchin representations 
    $\mathcal{H}(\Gamma) := \{\rho \in \mathcal{A}_\R(\Gamma) : \rho \text{in the same Anosov component as }\rho_0\}$
    \item Quasi-Hitchin representations
    \item $SL_3\C$-reducible representations
    \item Barbot representations
    \item Quasi-Barbot representations
\end{itemize}

%%%%%%%%%%%%%%%%%%%%%%%%%%%%%%
\section{Richardson varieties in $\F^3$}
TODO: This maybe belongs in an introduction / background section instead.

For any flag $x \in \F^3$, we define $Sch(x)$ the principle Schubert varieties associated with $x$. 
Let $x = (l, H)$ where $l \in \CP^2$ and $H \in (\CP^2)^*$ such that $l_x \subset H_x$.
We distinguish the elements of $Sch(x)$ using elements of the Weyl group $W(SL_3\C)$ as follows:
\begin{itemize}
    \item $\alpha_{Id}(x) = \{ (l, H) \} = \{x\}$,
    \item $\alpha_1(x) = \{ (l, H_l) : l \subset H_l \}$,
    \item $\alpha_2(x) = \{ (l_H, H) : l_H \subset H \}$,
    \item $\alpha_{12}(x) = \{ (k, H_l) : l \subset H_l \}$,
    \item $\alpha_{21}(x) = \{ (l_H, G) : l_H \subset H \}$, and
    \item $\alpha_{121}(x) = \alpha_{212}(x) = \F^3$.
\end{itemize}
The principle Schubert varieties are closed, and have complex dimensions $(0,1,1,2,2,3)$, respectively.

We call a pair of flags $x,y \in \F^3$ \emph{transverse} if they are in general position, 
meaning that for any proper, principle Schubert variety $\alpha$, 
then $x \not \in \alpha(y)$ and $y \not \in \alpha(x)$.

Let $x,y \in \F^3$ transverse. 
Define $R^0_{x,y}$ to be their irreducible 0-dimensional Richardson varieties, 
which are the intersections between principle Schubert varieties $\alpha(x) \cap \beta(y)$ 
where $dim(\alpha) + dim(\beta) = dim(\F^3)$.
There are exactly six such varieties $\{x, x_y, x^y, y_x, y^x, y\}$, where:
    \begin{itemize}
        \item $x := \{x\} \cap \alpha_{121}(y)$,
        \item $x_y := \alpha_1(x) \cap \alpha_{12}(y)$,
        \item $x^y := \alpha_2(x) \cap \alpha_{21}(y)$,
        \item $y_x := \alpha_{12}(x) \cap \alpha_1(y)$,
        \item $y^x := \alpha_{21}(x) \cap \alpha_2(y)$, and
        \item $y := \alpha_{121}(x) \cap \{y\}$.
    \end{itemize}
Given two tranverse flags, the Richardson 0-varieties can be computed explicitly:
\begin{example}\label{eg:r0}
    Choose $x$ and $y$ are the attracting and repelling flags of 
    $\rho(g) = diag(\lambda, \mu, \nu)$ where $|\lambda| > |\mu| > |\nu|$.
    Then the six Richardson 0-varieties in $R^0_{x,y}$ are
    $x = (e_1, e_1 + e_2)$, 
    $x_y = (e_1, e_1+e_3)$, 
    $x^y = (e_2, e_1+e_2)$, 
    $y_x = (e_3, e_1+e_3)$,
    $y^x = (e_2, e_2+e_3)$, and
    $y = (e_3, e_2 + e_3)$.  
\end{example}

Let $x,y \in \F^3$ transverse, and
define $R^1_{x,y}$ to be irreducible 1-dimensionsal Richardson varieties.
Similarly, these are intersections of principle Schubert varieties 
$\alpha(x) \cap \beta(y)$ where $dim(\alpha) + dim(\beta) = dim(\F^3) + 1$. 
There are six such 1-varieties, each homeomorphic to $\CP^1$, 
which we distinguish using pairs of elements of $R^0_{x,y}$.
\begin{itemize}
    \item $span(x, x_y) := \alpha_1(x) \cap \alpha_{121}(y)$,
    \item $span(x, x^y) := \alpha_2(x) \cap \alpha_{121}(y)$,
    \item $span(x_y,y_x) := \alpha_{12}(x) \cap \alpha_{12}(y)$,
    \item $span(x^y,y^x) := \alpha_{21}(x) \cap \alpha_{21}(y)$,
    \item $span(y_x, y) := \alpha_{121}(x) \cap \alpha_1(y)$, and
    \item $span(y^x, y) := \alpha_{121}(x) \cap \alpha_2(y)$.
\end{itemize}

The 1-varieties are not pairwise disjoint, e.g. $\{x\} = span(x,x_y) \cap span(x,x^y)$.
The points of intersection between 1-dimensional Richardson varieties 
are the 0-dimensional Richardson varieties. 

\begin{remark} 
    If $x,y \in \F^3$ transverse and 
    $\alpha,\beta$ principle Schubert varieties such that $dim(\alpha)+dim(\beta) < dim(\F^3)$, 
    then $\alpha(x) \cap \beta(y) = \null$.
\end{remark}

%%%%%%%%%%%%%%%%%%%%%%%%%%%%%%%%%%%%%%
\section{Richardson varieties in $\F^{n}$ for $n\geq 4$}

Conjectural picture for Richardson varieties in higher dimensions:

I suspect that topological conjugacies of Borel Anosov representations will preserve 
Richardson varieties up to $R^{n-2}_{x,y}$ organized into a cell decomposition with data an 
$(n-1)$-valent graph where vertices correspond to 0-varieties, edges correspond to 1-varieties,
fully connected triples correspond to 2-varieties, etc. 
Matching/switching Richardson varieties will mean sending Richardson varieties to their dual.

I think balanced ideals $I$ correspond to subgraphs of the Richardson variety graph 
containing exactly half of the vertices. In other words, $Th^I(x) \cap Th^I(y) = \null$, and
\begin{align}
    (Th^I(x) \cap R^0_{x,y}) \cup (Th^I(y) \cap R^0_{x,y}) = R^0_{x,y}.
\end{align}
However, topological conjugacies should not depend on the choice of balanced ideal, 
so then they determine diffeomorphisms between all types of quotient manifolds. 
In other words, if $F:\F^n \rightarrow \F^n$ a topological conjugacy between $\rho, \rho'$, then:
\begin{align}
    F(M^I_\rho) = M^I_{\rho'} \text{ for all balanced ideals } I \subset W(SL_n\C).
\end{align}
Most importantly, I believe the conclusions of Theorem \ref{thm:hyperconvex-conjugate} will apply in higher rank:
\begin{conjecture}
    For $\rho,\rho':\Gamma \rightarrow SL_n\C$ Borel Anosov and topologically conjugate. 
    Then $\rho$ hyperconvex if and only if $\rho'$ hyperconvex. 
\end{conjecture}